\chapter{Einleitung}
% \abschnittsautor controls who appears on the right of the page header. It
% works per section too, not just per chapter -- call it again after any
% \section that a different team member wrote. ALWAYS call it directly after
% the heading, otherwise it takes effect one page too early.
\abschnittsautor{\alleautoren}

\section{Motivation}

Die Einleitung führt in das Thema ein: Warum ist die Aufgabenstellung
relevant, in welchem Umfeld tritt sie auf, und was passiert, wenn sie
ungelöst bleibt? Belege gehören hier bereits mit Quellenangabe
dazu~\cite{beispiel-buch}.

\section{Zielsetzung}

Beschreibt das Ziel so, dass am Ende der Arbeit überprüfbar ist, ob ihr es
erreicht habt. Trennt dabei funktionale von nichtfunktionalen Anforderungen.

Ein Absatz wird durch eine Leerzeile im Quelltext erzeugt. Einen Zeilenumbruch
innerhalb eines Absatzes erzwingt ihr mit \verb|\\| -- braucht ihr aber fast nie.
